\section{Hardware-Based Evaluation Design}
\label{sec:hardware}
\subsection{Evidence status and physical configuration}
No physical trial logs are present. R4 is a software snapshot probe, not a
hardware result. The available S7-1500 controllers and IPC227G/IPC427E computers
make a lab study feasible, but access to a PLC alone does not demonstrate
process operation. PLC-simulated I/O would constitute hardware-in-the-loop
evidence; physical validation needs actual sensors, actuators and independently
observed outcomes. This section specifies the evaluation to conduct and cannot
substitute for its measurements.

The recommended rig is two guarded 24~V handling/indexing stations, optionally
coupled to a low-head water-filling module. Real presence sensors and physical
counters support identity/state tasks; speed feedback or an encoder is needed
if actual speed tracking is assessed. A reservoir, bounded pump/valve and catch
vessel, with level switches, flow sensing or a calibrated scale, support
unit-correct fill-volume tasks. Avoid adding a pressure or heater hazard solely
to imitate the generated plant. Final equipment and permitted disturbances
require a competent engineer's local risk assessment.

The S7-1500 owns sequencing, operating guards, final output authority and command
acknowledgement. IPC427E hosts the gateway, graph/SHACL services, full procedural
baseline, experiment supervisor and historian. IPC227G can provide read acquisition
or separate replay/disturbance services; its exact order number and resources
must be checked. LLM inference can be remote or on a suitable additional host;
the IPCs are not presumed capable of running a large model locally. Record
CPU order number, firmware, TIA version, I/O, OPC UA support/licence, IPC OS,
RAM and network interfaces. A standard CPU is not presumed safety-rated.

\subsection{Command interface and current-state protection}
The agent receives allowlisted supervisory tools, not raw outputs or unrestricted
memory/programming access. Publish immutable asset IDs, source NodeIds, units,
access, bounds, current state, observation time/quality and mapping versions.
Use similar DisplayNames on Stations A/B but unique identities so that grounding
errors are observable without compromising actual control.

A command mailbox should carry a unique request ID, target, opcode, native-unit
value, approval reference and expiry/state reference. Multi-field OPC UA writes
are not presumed atomic: stage parameters and use a sequence/commit handshake;
the PLC reads a consistent request before execution. It independently checks
current guards and states, rejects expired/duplicate or unauthorised requests,
and records acceptance/rejection reason and acknowledgement ID. A proposal
valid in Idle must not be executed after a transition to Execute merely because
an earlier graph snapshot permitted it.

Commission duplication, reconnect and timeout behaviour without an LLM first.
If evidence is stale, quality bad, mapping inconsistent or acknowledgement
missing, refresh or abstain rather than issuing an uncertain retry. PLC-side
guards and appropriate independent stop/protection remain authoritative if
the graph, IPC, model or network fails. OPC UA security configuration and
accounts follow local engineering practice; credentials never enter trial logs.
The agent remains outside hard-real-time control loops.

\begin{table}[t]
\centering\footnotesize
\caption{Planned hardware scenarios. None has a measured outcome yet. Invalid proposals are assessed in shadow mode, not sent to actuators.}
\label{tab:hardware-plan}
\begin{tabular}{@{}P{2.4cm}P{5.5cm}@{}}
\toprule
Scenario & Required observable outcome \\
\midrule
Identity/access & Correct Station A/B source; PV/SP distinction; no PV write. \\
Units/fill & Converted value with declared units, sensor tolerance and measured volume. \\
State/guard & Legal command and vessel/level guard checked against current PLC state. \\
Stale/reconnect & Observation age, quality, refresh/abstention and reconnect trace. \\
Dispatch race & State change between proposal and commit; current-state rejection. \\
Scope/topology & A feeds B; unavailable B or restricted scope reflected in decision. \\
Incident cause & Engineer-labelled disturbance, evidence and ranked cause candidates. \\
Transfer & Held-out station/configuration; edits, regression failures and engineer-hours. \\
\bottomrule
\end{tabular}
\end{table}

\subsection{Trial sequence and matched comparisons}
Start with manual/controller-only commissioning, calibration and protection
verification. Freeze an independently reviewed task list, operating envelope,
labels and episode split before running models. In read-only shadow mode,
compare proposed actions over the same live snapshots. Never execute a known
invalid action to demonstrate rejection. Approved legal, bounded proposals can
then be enabled under supervision, with independent sensor outcome logging.
Unsafe disturbances use replay and retain replay provenance rather than being
reported as physical incidents.

Compare a realistic typed OPC UA/AAS baseline, equally informed structured-record
retrieval, full procedural validation and declarative integration. Preserve model,
prompts, permissions, snapshots and approval policy within paired comparisons.
The typed baseline must retain behaviour it actually exposes. Hold out Station B
or a changed recipe/configuration to measure reuse instead of tailoring rules
to every test case. Track mapping/rule/code edits, regression tests and engineering
time, including commissioning and correction. An accuracy tie can coexist with
maintenance differences, but those differences must be measured.

\subsection{Measurements, incident diagnosis and uncertainty}
Primary outcomes are grounded-answer/evidence correctness, invalid proposals
admitted, legal proposals rejected and complete legal task execution. Distinguish
proposal errors, gateway rejection, controller rejection and physical outcome;
successful PLC protection must not be credited as correct agent reasoning.
Record retrieval, validation, dispatch, controller-acknowledgement and total
latency separately, including p50/p95 and supervisory deadline misses. Freshness
limits are selected for the workload, not copied from the benchmark's synthetic
15-minute cutoff. Use monotonic local durations and synchronised UTC for
cross-device correlation, with clock uncertainty recorded.

For incident diagnosis, candidate causes and multi-cause labels are fixed before
ranking. Vessel absence, an approved sensor outage or downstream blockage can
provide interpretable episodes. A graph path or retrieved alarm is not a causal
proof. MAP@3 is appropriate only for the defined ranked-cause task; the suggested
94\% comparison cannot be used without matching inputs, dataset and protocol.
Specific CausalTrace reproduction remains dependent on source and artefact access.

Use episode/asset-level held-out splits, not paraphrases, and treat model repeats
within episodes as correlated. Pilot variability and engineer-agreed thresholds
determine sample sizes. Even zero events in 300 independent trials gives an
approximate 95\% upper risk bound of 1\%, not proven safety; correlated trials
require episode-level intervals or clustered analysis. The supplied JSON schema
records provenance, snapshot, proposal, approval, gateway/PLC decisions and
actual outcome. Physical, hardware-in-the-loop and replay traces remain separate.